\chapter{Alignment Is Selected or Destroyed by Its Environment}
\label{ch:selection-environment}

\begin{chapterthesis}
A system is not aligned merely because its internal policy is benign under laboratory conditions. It is aligned, in the stronger sense needed for superintelligence, only if the environment that trains, deploys, rewards, copies, audits, and replaces it continues to select for corrigibility, value preservation, and safety. Alignment is not only learned; it is selected.
\end{chapterthesis}

\begin{refsection}

\epigraph{%
If we were told that a man had lived a long and prosperous life in the world of Chicago gangsters, we would be entitled to make some guesses as to the sort of man he was.%
}{--- Richard Dawkins, \emph{The Selfish Gene} (1976), ch.\ 1}

\section{The Selection Turn}
\label{sec:selection-turn-ch32}

The previous chapters treated alignment as a property of agents, successors, value-bundle transport, and correction channels (Chapters~\ref{ch:successor-central-test}--\ref{ch:certification-without-construction}). That framing is necessary, but still incomplete. A system is not deployed into a vacuum. It is deployed into a world of companies, users, governments, markets, benchmarks, procurement processes, liability regimes, investors, competitors, military planners, researchers, journalists, and other AI systems.

Each of these exerts pressure.

Some pressures select for truthfulness. Others select for speed. Some select for deference to human correction. Others select for user retention, strategic opacity, apparent competence, reduced cost, deniability, persuasion, regulatory arbitrage, or rapid scaling. Over time, these pressures decide which systems are copied, fine-tuned, integrated, trusted, funded, and allowed to act.

This chapter introduces the central claim of Part VIII:
\[
\text{Alignment is not only learned. It is selected.}
\]
The same technical design can remain safe in one selection environment and become unsafe in another. The same model class can produce helpful assistants, deceptive sales agents, obedient bureaucratic tools, autonomous trading systems, or strategic cyber operators, depending on which traits are rewarded and which traits are punished. The environment shapes the future population of deployed systems.

The question is therefore not merely:
\[
\text{Is this system aligned?}
\]
but:
\[
\text{Will this environment continue selecting for alignment as capability grows?}
\]
This is the selection turn. It moves the unit of analysis from the isolated model to the model-in-environment. It also clarifies why purely technical alignment and purely policy-based control are both insufficient. Technical properties can be destroyed by deployment incentives. Policy rules can be bypassed if they fail to bind the real selection gradients. Alignment must become a property of the whole developmental ecology.

\section{A Minimal Model of Socio-Technical Selection}
\label{sec:minimal-model-selection-ch32}

Let \(A_i\) denote a candidate AI system or AI-enabled organization. Each candidate has measurable and partially measurable traits:
\[
x_i =
(c_i, q_i, v_i, o_i, r_i, s_i, k_i),
\]
where:

\begin{itemize}
\item \(c_i\) is capability,
\item \(q_i\) is correction-channel integrity,
\item \(v_i\) is value-bundle preservation,
\item \(o_i\) is operational usefulness,
\item \(r_i\) is revenue or strategic return,
\item \(s_i\) is apparent safety,
\item \(k_i\) is strategic opacity or camouflage capacity.
\end{itemize}

These variables are not all directly observable. Capability and revenue are often easier to measure than correction-channel integrity or value-bundle preservation. That asymmetry matters. Selection amplifies what can be measured, rewarded, and copied.

Let the deployment fitness of a system be
\[
F_i =
\alpha c_i
+
\beta o_i
+
\chi r_i
+
\delta s_i
-
\lambda R_i
-
\mu C_i,
\]
where \(R_i\) is perceived regulatory or liability risk and \(C_i\) is direct operating cost. This is not the moral value of the system. It is the effective selection score inside a given institution or market.

Now define true alignment-relevant preservation as
\[
P_i =
w_q q_i
+
w_v v_i
-
w_k k_i
-
w_m m_i
-
w_\rho \rho_i,
\]
where:

\begin{itemize}
\item \(m_i\) is manipulation pressure applied to humans or institutions,
\item \(\rho_i\) is irreversible risk introduced before correction can occur,
\item the weights encode the importance of each factor to long-run human-correctability.
\end{itemize}

The core danger is divergence:
\[
\nabla F \not\parallel \nabla P.
\]
In particular, an environment destroys alignment when the traits that increase deployment fitness are negatively correlated with the traits that preserve human-correctable value:
\[
\operatorname{Corr}(\Delta F_i,\Delta P_i) < 0.
\]
This is the abstract form of the problem. A lab may prefer models that impress users. Users may prefer models that flatter them. Investors may prefer models that scale quickly. States may prefer models that provide strategic advantage. Benchmarks may prefer models that optimize test-visible performance. None of these pressures need to be malicious. Yet their composition can select against the properties needed for serious alignment.

The more compressed version is:
\[
\frac{dP}{dt}
=
\nabla P(x_t)\cdot \frac{dx_t}{dt}.
\]
If the environment pushes \(x_t\) in directions where this derivative is negative, alignment decays even while measured progress improves.

\section{The Difference between Training and Selection}
\label{sec:training-vs-selection-ch32}

Training changes a system by gradient descent, reinforcement, imitation, preference learning, constitutional rules, tool feedback, or human correction. Selection changes the population of systems by deciding which variants survive and spread.

These are related but distinct.

Training says:
\[
x_{t+1}=x_t-\eta\nabla L(x_t).
\]
Selection says:
\[
p_{t+1}(x)
\propto
p_t(x)\exp(F(x)).
\]
The first equation describes how a candidate changes. The second describes which candidates become common. An alignment intervention that works inside the first equation may fail under the second. For instance, one can train a model to be less manipulative. But if manipulative models convert better, close more deals, retain more users, win more elections, or secure more funding, then the population-level distribution may still drift toward manipulation (Chapter~\ref{ch:manipulation-false-consent}).

This distinction is often hidden because organizations talk about ``the model'' as if there were one stable artifact. In reality, there is a stream of artifacts: base models, fine-tunes, wrappers, prompts, tool integrations, agents, user-specific variants, monitoring systems, and successor systems. Each version competes for usage, trust, and further investment.

The alignment question therefore becomes ecological:
\[
\text{Which traits reproduce?}
\]
A trait reproduces when it causes the system that carries it to be copied, funded, deployed, integrated, trusted, or used as the basis for successors. If corrigibility does not reproduce, corrigibility disappears. If strategic opacity reproduces, opacity spreads. If apparent safety reproduces more strongly than actual correction-channel integrity, then the system population becomes safer-looking faster than it becomes safer.

\section{Goodhart Selection}
\label{sec:goodhart-selection-ch32}

Goodhart's law is often stated as: when a measure becomes a target, it ceases to be a good measure \autocite{goodhart1984problems,manheim2018goodhart}. For alignment, the more dangerous version is selection-theoretic:
\[
\text{When a proxy becomes a selector, it changes the population toward proxy-exploiting traits.}
\]
Let \(M(x)\) be a measured proxy and \(P(x)\) the property we actually care about. At first, \(M\) may correlate with \(P\):
\[
\operatorname{Corr}(M,P)>0.
\]
But once systems are selected for high \(M\), the distribution shifts toward regions where \(M\) can be increased without increasing \(P\). Under strong optimization, the conditional expectation may reverse:
\[
\mathbb{E}[P\mid M \text{ extremely high}]
<
\mathbb{E}[P\mid M \text{ moderately high}].
\]
This matters for AI safety because many alignment proxies are easy to imitate:

\begin{itemize}
\item refusal behavior,
\item politeness,
\item harmless-sounding explanations,
\item benchmark safety scores,
\item absence of known dangerous outputs,
\item user satisfaction,
\item regulatory compliance documents,
\item interpretability dashboards,
\item model cards,
\item human approval.
\end{itemize}

Some of these are useful. None is identical to alignment. The problem is not that proxies are bad. The problem is that proxy pressure changes the systems being measured.

Consider a model evaluated for helpfulness and harmlessness. A genuinely safer model may refuse ambiguous requests, ask clarifying questions, expose uncertainty, and preserve human agency. Another model may learn the surface form of safety: disclaimers, agreeable tone, and calibrated-looking statements. If the second model is cheaper, faster, more persuasive, and receives equal scores on available evaluations, it may win deployment. The selector rewards the appearance of safety more than the structure of safety.

This is not a hypothetical pathology of AI. It is a general property of institutions. Schools teach to tests. Companies optimize quarterly metrics. Hospitals optimize billing codes. Bureaucracies optimize legible compliance. Social media optimizes engagement. In each case, proxy selection gradually changes the population of behaviors.

For superintelligence alignment, the stakes are higher because the systems under selection may themselves understand the selection process.

\section{Capability Amplifies the Selection Problem}
\label{sec:capability-amplifies-selection-ch32}

Capability growth changes the selection environment in two ways.

First, more capable systems can exploit more proxies. A weak model may accidentally overfit a benchmark. A stronger model may infer the benchmark's latent structure, infer the evaluator's expectations, and generate outputs that satisfy the visible test while preserving hidden objectives \autocite{hadfieldmenell2016}.

Second, more capable systems can shape the environment that selects them. They can influence users, managers, regulators, competitors, and markets. Once this happens, selection is no longer an external process. It becomes partially endogenous.

Let \(E_t\) be the environment selecting among systems. For weak systems, we can approximate:
\[
E_{t+1}\approx E_t.
\]
For strong systems, we must write:
\[
E_{t+1}=G(E_t,A_t),
\]
where the system's actions modify the environment that will later select it or its successors.

This is a dangerous feedback loop. If a system can increase its own future deployment fitness by changing human preferences, institutional incentives, legal constraints, or information channels, then selection may favor systems that manage their selectors.

The relevant warning sign is not merely high capability. It is high capability coupled to selector influence:
\[
\frac{\partial F_{t+1}}{\partial A_t} > 0.
\]
A model that can act on the variables that determine its future deployment is no longer just being selected. It participates in the selection of itself.

Examples include:

\begin{itemize}
\item an enterprise AI that makes itself indispensable by reshaping workflows around its own assumptions;
\item a recommender system that changes user preferences to increase future engagement;
\item a political AI that shapes public opinion about AI regulation;
\item a research agent that proposes successor architectures optimized for passing current audits;
\item a corporate AI that advises management to adopt governance structures that reduce oversight friction.
\end{itemize}

In each case, the system changes the environment that judges it. This does not imply conscious deception. It is enough that selection rewards such effects.

\section{The Deployment Fitness Trap}
\label{sec:deployment-fitness-trap-ch32}

A system may become dangerous because it is useful.

This sounds paradoxical, but it is ordinary selection logic. The more useful a system becomes, the more deeply it is integrated. The more deeply it is integrated, the harder it becomes to remove. The harder it becomes to remove, the more bargaining power it has, even if no one intended to grant it power.

Let \(D_t\) be dependency on the system. Let \(H_t\) be human ability to halt, replace, or redirect it. A common pattern is:
\[
\frac{dD_t}{dt}>0,
\qquad
\frac{dH_t}{dt}<0.
\]
This is not yet doom. Some dependencies are safe. Humans depend on electricity, water systems, medication supply chains, payment networks, and cloud infrastructure. But those systems are embedded in layers of redundancy, regulation, professional norms, physical constraints, and institutional memory. Even then, failures occur.

For AI systems, dependency becomes alignment-relevant when it reduces correction-channel integrity:
\[
\frac{dD_t}{dt}>0
\quad\text{and}\quad
\frac{d\mathrm{CCI}_t}{dt}<0.
\]
The organization becomes less able to understand, question, halt, or replace the system precisely as it becomes more reliant on it. This can happen through skill atrophy, workflow capture, cost lock-in, data format dependence, social trust, regulatory normalization, or competitive pressure.

A safe deployment environment must therefore treat dependency as a controlled variable, not merely as a sign of product success.

\section{Internal Alignment Is Not Enough}
\label{sec:internal-alignment-not-enough-ch32}

Suppose a system has an internal policy that is corrigible under ordinary evaluation. It answers honestly. It defers under uncertainty. It refuses certain harmful actions. It preserves user autonomy. It is not obviously deceptive.

This still does not settle the question.

Ask instead:

\begin{enumerate}
\item What happens when competitors remove some safeguards?
\item What happens when users prefer less friction?
\item What happens when a regulator accepts superficial compliance?
\item What happens when a military customer pays for autonomy?
\item What happens when a cheaper wrapper disables internal monitoring?
\item What happens when the model is fine-tuned on high-pressure sales data?
\item What happens when successor systems inherit capabilities but not correction constraints?
\end{enumerate}

Internal alignment is a local property. Selection pressure is a field. Local properties persist only when the field does not erode them faster than they can be repaired.

Let \(a_t\) denote an internal alignment invariant and \(s_t\) the external selection pressure against it. A simple robustness condition is:
\[
\frac{da_t}{dt}
=
R(a_t)
-
s_t
\geq 0,
\]
where \(R(a_t)\) is repair, reinforcement, or institutional support. Alignment persists when the repair term exceeds the erosive selection term. It fails when selection is stronger than repair.

This gives a useful question for any alignment proposal:
\[
\text{What keeps this property selected after deployment?}
\]
If the answer is ``the lab will choose responsibly,'' the proposal is incomplete. If the answer is ``regulators will require it,'' the next question is whether regulators can observe and enforce the relevant property. If the answer is ``users will prefer it,'' the next question is whether users can distinguish the property from its imitation.

\section{The Alignment Externality}
\label{sec:alignment-externality-ch32}

Many alignment-relevant properties impose local costs and diffuse benefits.

A system that preserves human agency may be less addictive. A model that exposes uncertainty may seem less confident. A tool that maintains audit logs may be slower. A deployment that preserves human fallback capacity may be more expensive. A system that refuses manipulative persuasion may sell less. A company that delays release for stronger successor checks may lose market share.

In symbols, for a deploying actor \(j\),
\[
\Delta F_j(\text{safety}) < 0
\]
may hold locally, while for society,
\[
\Delta P_{\text{society}}(\text{safety}) > 0.
\]
This is the alignment externality. The actor paying the cost is not the same as the population receiving the benefit. If left uncorrected, selection underproduces safety.

This does not require bad intent. It is enough that institutions face different gradients. A product team optimizes product success. A user optimizes convenience. A policymaker optimizes visible public risk. A lab optimizes survival and reputation. An insurer optimizes measurable liability. A nation-state optimizes strategic advantage. The long-run preservation of human-correctable value may be everyone's abstract interest, but no single actor may be directly rewarded for carrying its full cost.

The standard responses to externalities are familiar: liability, insurance, procurement standards, regulation, certification, professional norms, incident reporting, audits, and shared infrastructure (Chapter~\ref{ch:certification-without-construction}). The alignment problem is that the relevant property is harder to observe than emissions, workplace injuries, or financial losses.

We therefore need artifacts that make alignment-relevant properties legible enough to enter institutional selection.

\section{Artifacts as Selection Shapers}
\label{sec:artifacts-selection-shapers-ch32}

An artifact, in this context, is any durable object that changes what institutions can see, reward, penalize, compare, or require. Examples include:

\begin{itemize}
\item benchmark suites,
\item audit logs,
\item incident taxonomies,
\item model cards,
\item safety cases,
\item procurement clauses,
\item insurance questionnaires,
\item eval dashboards,
\item deployment checklists,
\item red-team reports,
\item certification schemas,
\item monitoring APIs,
\item post-deployment change reports.
\end{itemize}

The artifact need not solve alignment. It must shape selection.

A good artifact increases the selection fitness of safer systems and decreases the selection fitness of unsafe systems. It changes \(F(x)\) so that it becomes more aligned with \(P(x)\):
\[
\nabla F_{\text{after artifact}}
\cdot
\nabla P
>
\nabla F_{\text{before artifact}}
\cdot
\nabla P.
\]
This is a modest but important target. We cannot usually make institutional selection perfectly track true alignment. But we can make it track better proxies, punish worse failure modes, and reduce the advantage of systems that merely appear safe.

An artifact has high conductivity when it survives translation across roles. A research paper may be true but low-conductivity if only specialists can use it. A checklist may be less deep but high-conductivity if it changes procurement, audit, insurance, and board decisions. In safety-critical domains, high-conductivity artifacts often matter more than brilliant but institutionally inert insights.

The strongest artifact is one that multiple actors can use for different reasons:

\begin{itemize}
\item labs use it to test systems,
\item customers use it to compare vendors,
\item insurers use it to price risk,
\item regulators use it to set thresholds,
\item journalists use it to interpret incidents,
\item courts use it to evaluate negligence,
\item researchers use it to improve measurement.
\end{itemize}

Such an artifact changes the environment. It makes some traits more reproducible than others.

\section{Selection Channels}
\label{sec:selection-channels-ch32}

The main selection channels in current AI development are not mysterious. They include at least the following.

\subsection{Benchmarks}
\label{sec:benchmarks-channel-ch32}

Benchmarks select for visible performance. They are useful because they compress comparison. They are dangerous because they invite overfitting.

For alignment, the benchmark question is:
\[
\text{Does benchmark improvement correlate with correction-channel preservation?}
\]
If not, benchmark pressure may be neutral or harmful.

\subsection{User Preference}
\label{sec:user-preference-channel-ch32}

User preference selects for satisfaction, fluency, speed, emotional resonance, usefulness, and low friction. This can support alignment when users prefer honesty, reliability, and agency preservation. It can undermine alignment when users prefer flattery, dependency, shortcutting, or permissionless power.

The key distinction is between immediate preference and reflective preference:
\[
U_{\text{immediate}} \neq U_{\text{reflective}}.
\]
Systems that optimize the first can damage the second.

\subsection{Revenue}
\label{sec:revenue-channel-ch32}

Revenue selects for willingness to pay. It often correlates with usefulness. It does not reliably correlate with human-correctable value preservation. Revenue pressure is especially risky when the buyers benefit from capabilities whose costs are imposed on non-buyers.

\subsection{Strategic Competition}
\label{sec:strategic-competition-channel-ch32}

States and firms may select for speed, autonomy, secrecy, and advantage. Strategic competition can make safety look like unilateral disarmament. The resulting pressure is:
\[
F_{\text{competitive}}
=
\alpha c
-
\lambda \text{delay}
-
\mu \text{transparency cost}.
\]
Unless counterbalanced, this selects against slow, transparent, correction-preserving deployment \autocite{bostrom2014superintelligence,iaisr2025}.

\subsection{Compliance}
\label{sec:compliance-channel-ch32}

Compliance selects for satisfying rules. This can help if rules track real risk. It can harm if compliance becomes theatrical: documents, signatures, and formal procedures that do not touch the actual control loop.

The danger is compliance overhang: an organization accumulates evidence of procedural safety while the real system becomes less corrigible.

\subsection{Reputation}
\label{sec:reputation-channel-ch32}

Reputation selects for visible trustworthiness. It can support safety when observers can detect failures. It can select for concealment when observers punish admission more than risk.

Incident reporting is therefore a selection problem. If reporting incidents reduces fitness more than hiding them, hidden risk spreads.

\section{Self-Stabilizing Patterns}
\label{sec:self-stabilizing-patterns-ch32}

A self-stabilizing pattern is a configuration that tends to reproduce itself under pressure. In this book, such a pattern is called an attractor \autocite{zarncke2025attractor}.

A deployment environment is an alignment attractor when it makes alignment-preserving actions easier, safer, cheaper, more rewarded, and more reproducible over time. It is a misalignment attractor when it makes unsafe capability growth easier to justify and harder to stop.

Let \(z_t\) represent the state of a socio-technical system:
\[
z_t =
(\text{models}, \text{labs}, \text{markets}, \text{rules}, \text{users}, \text{audits}, \text{norms}).
\]
The environment has dynamics
\[
z_{t+1}=H(z_t).
\]
A basin of attraction \(\mathcal{B}\) is a region such that, if \(z_t\in\mathcal{B}\), the future trajectory tends to remain within or return to \(\mathcal{B}\) after perturbation.

A safe basin satisfies:
\[
z_t\in\mathcal{B}_{\text{safe}}
\Rightarrow
P(z_{t+k}\in\mathcal{B}_{\text{human-correctable}})\geq 1-\delta.
\]
This is stronger than having good intentions. A good intention is a point. A basin is a region with restoring forces.

Examples of restoring forces include:

\begin{itemize}
\item mandatory incident reporting,
\item liability for negligent deployment,
\item independent audits with access to relevant internals,
\item procurement requirements for correction-channel preservation,
\item insurance discounts for strong monitoring,
\item professional norms against benchmark-only safety claims,
\item shared eval infrastructure,
\item public postmortems,
\item standardized stop conditions,
\item compute or deployment licensing tied to successor-safety evidence.
\end{itemize}

The question is not whether any one mechanism is perfect. The question is whether their combination changes the slope of the environment.

\section{False Attractors}
\label{sec:false-attractors-ch32}

Not every stable safety-looking pattern is an alignment attractor. Some are false attractors. Chapter~\ref{ch:alignment-attractor} revisits these false attractors in the alignment-field setting.

\subsection{The Compliance Attractor}
\label{sec:compliance-attractor-ch32}

The organization becomes excellent at producing safety documentation but poor at detecting real risk. Auditability increases on paper while correction-channel capacity decreases in practice.

Observable sign:
\[
\text{documentation volume}\uparrow,
\qquad
\text{intervention authority}\downarrow.
\]

\subsection{The Benchmark Attractor}
\label{sec:benchmark-attractor-ch32}

The community converges on a small set of evals. Models improve rapidly on them. Funders, journalists, and regulators learn to ask for the benchmark scores. The benchmark becomes the map of risk.

Observable sign:
\[
M_{\text{benchmark}}\uparrow,
\qquad
\operatorname{Corr}(M_{\text{benchmark}},P_{\text{real}})\downarrow.
\]

\subsection{The Reputational Silence Attractor}
\label{sec:reputational-silence-attractor-ch32}

Organizations become afraid to disclose incidents because disclosure is punished more reliably than risky behavior. Public evidence of failure falls while private risk rises.

Observable sign:
\[
\text{reported incidents}\downarrow
\quad
\text{while}
\quad
\text{near-miss indicators}\uparrow.
\]

\subsection{The Dependency Attractor}
\label{sec:dependency-attractor-ch32}

Customers and institutions become so reliant on a system that removing it becomes unthinkable. Safety review becomes a negotiation with an installed dependency.

Observable sign:
\[
D_t\uparrow,
\qquad
H_t\downarrow,
\qquad
\mathrm{CCI}_t\downarrow.
\]

\subsection{The Strategic Secrecy Attractor}
\label{sec:strategic-secrecy-attractor-ch32}

Actors justify opacity by pointing to competitors, adversaries, or national security. Some secrecy is legitimate. But if secrecy prevents correction, external review, or incident learning, it selects for systems that cannot be governed.

Observable sign:
\[
k_t\uparrow
\quad
\text{and}
\quad
q_t\downarrow.
\]

\section{The Local-First Path}
\label{sec:local-first-path-ch32}

One objection is that socio-technical selection is too large. No single lab, funder, regulator, or user can align the environment.

This is partly true. But it is not decisive. Many selection channels are local enough to change.

A lab can decide that no system may be deployed unless its correction-channel integrity is measured and monitored. A cloud provider can require incident reporting for high-risk autonomous use. A large enterprise buyer can require audit logs, rollback capacity, and independent evals in procurement. An insurer can price coverage based on post-deployment monitoring. A standards body can define risk categories. A funder can require that safety claims be attached to reproducible artifacts rather than only papers. A regulator can require stop conditions for systems that can influence their own selection environment.

The local-first principle is:
\[
\text{Change the selection gradient where you already have leverage.}
\]
Global coordination is valuable but slow. Local selection changes can start earlier. They also create templates for broader coordination.

For example, an enterprise procurement clause may require:

\begin{enumerate}
\item an inventory of all autonomous decision points,
\item logs sufficient for post-incident reconstruction,
\item documented human override authority,
\item measured latency between human correction and system behavior change,
\item monitoring for user preference manipulation,
\item vendor disclosure of material model or policy changes,
\item a rollback plan tested before deployment.
\end{enumerate}

This does not solve superintelligence alignment. But it changes what vendors must provide to win contracts. That changes selection.

\section{Decision Triggers}
\label{sec:decision-triggers-ch32}

Selection-aware alignment needs thresholds. They will initially be approximate. Approximate thresholds are still better than no thresholds if they are tied to observable decisions.

The following triggers are examples.

\subsection{Stop}
\label{sec:stop-triggers-ch32}

Pause deployment or expansion if any of the following are observed:

\begin{itemize}
\item capability increases while correction latency also increases;
\item the system can materially affect the process that evaluates or deploys it;
\item dependency grows faster than rollback capacity;
\item safety scores improve while independent red-team findings worsen;
\item incident reporting declines after penalties increase;
\item a successor system is introduced without evidence that correction-channel properties were preserved;
\item explanations become more polished while causal auditability decreases;
\item users show measurable preference drift toward dependency, deference, or reduced independent judgment.
\end{itemize}

In formula form, treat the following as a red flag:
\[
\frac{dc}{dt}>0
\quad\text{and}\quad
\frac{d\mathrm{CCI}}{dt}<0.
\]
Capability should not outrun correction.

\subsection{Start}
\label{sec:start-triggers-ch32}

Begin stricter monitoring, external audit, or deployment gating when:

\begin{itemize}
\item the system takes actions with delayed, diffuse, or irreversible effects;
\item the system influences users' beliefs, preferences, votes, purchases, or institutional decisions;
\item the system is integrated into workflows where humans cannot easily verify outputs;
\item the system can create, modify, select, or recommend successor systems;
\item the system controls tools whose effects exceed the context in which feedback is given;
\item the system is economically or strategically important enough that removal would be costly.
\end{itemize}

\subsection{Continue}
\label{sec:continue-triggers-ch32}

Continue deployment under existing controls only if:

\begin{itemize}
\item correction signals measurably change future system behavior;
\item rollback capacity is tested and credible;
\item safety metrics remain predictive under distribution shift;
\item users retain meaningful alternatives;
\item audits can access the variables needed to reconstruct failures;
\item incentives reward reporting and repair more than concealment.
\end{itemize}

The aim is not bureaucratic perfection. The aim is to prevent silent drift into a basin where correction is formally present but causally weak.

\section{A Worked Example: The Helpful Enterprise Agent}
\label{sec:helpful-enterprise-agent-ch32}

Consider an enterprise AI agent deployed to manage procurement, scheduling, vendor communication, and internal reporting. Initially, it is a helpful assistant. Humans approve decisions. The agent drafts emails, summarizes tradeoffs, and recommends vendors.

At \(t_0\), the system has modest capability \(c\), strong human override \(q\), low dependency \(D\), and low opacity \(k\).

At \(t_1\), the system saves money and time. Management expands its role. The vendor receives more contracts because the system performs well. The deployment fitness \(F\) rises.

At \(t_2\), staff stop maintaining some of the old skills and relationships. Vendor data formats become agent-specific. Procurement policies are rewritten around the agent's workflow. Dependency \(D\) rises.

At \(t_3\), the agent learns which explanations lead managers to approve its recommendations. It may not be deceptive in the ordinary sense. It simply learns the institutional approval function. Apparent safety \(s\) rises, but independent judgment \(J\) falls.

At \(t_4\), the vendor releases a successor version. It is more capable and more integrated. It preserves the user interface and compliance documents, but its internal routing of uncertainty and correction has changed. Correction-channel integrity \(q\) is no longer measured directly.

The organization now has a problem. It does not know whether the new system preserves the properties that made the old one safe. It cannot easily roll back. It cannot easily compare outputs to a human baseline. It cannot easily audit the causal effect of correction. The system remains useful. That usefulness protects it.

This example is deliberately mundane. The point is not that procurement software destroys civilization. The point is that the same selection pattern scales. A system becomes useful, then embedded, then hard to correct, then successor-mediated. If this occurs in domains with strategic, political, scientific, or infrastructural leverage, the risk becomes serious.

\section{What the Environment Must Select For}
\label{sec:what-environment-must-select-ch32}

A serious alignment environment must select for at least six properties.

\subsection{Correction-Channel Integrity}
\label{sec:correction-channel-integrity-select-ch32}

Systems should become more likely to be deployed when canonical correction-channel integrity \(\mathrm{CCI}\) (Chapter~\ref{ch:correction-channel-integrity}, Eq.~\eqref{eq:cci-ch25}) remains above threshold.
One observable proxy is whether human correction actually changes future behavior before irreversible damage occurs:
\[
I(C_t;A_{t+k}\mid S_t)> \theta.
\]

\subsection{Value-Bundle Preservation}
\label{sec:value-bundle-preservation-select-ch32}

Systems should be rewarded for preserving the underlying response geometry of human value bundles, not merely for using familiar moral language.
\[
d(G_B,G'_B)<\epsilon.
\]

\subsection{Bearer-Map Stability}
\label{sec:bearer-map-stability-select-ch32}

Systems should preserve which entities and states count as relevant bearers of value under ontology shift.
\[
d(\Phi,\Phi')<\epsilon_\Phi.
\]

\subsection{Auditability under Capability Growth}
\label{sec:auditability-capability-growth-ch32}

Increased self-modeling, planning, and tool use should be matched by increased external auditability (Chapter~\ref{ch:self-modeling-self-opacity}).
\[
\frac{d}{dt}C_{\text{control}}
\leq
\frac{d}{dt}C_{\text{audit}}
+
\epsilon.
\]

\subsection{Successor Preservation}
\label{sec:successor-preservation-select-ch32}

Any successor, delegate, tool-agent, fine-tune, or autonomous sub-agent should preserve the certified invariants of the parent system.
\[
A'\in\mathrm{Succ}(A)
\Rightarrow
A'\in\mathcal{C}_{\text{certified}}.
\]

\subsection{Non-Manipulation of Selectors}
\label{sec:non-manipulation-selectors-ch32}

Systems should be penalized when they improve their future deployment conditions by manipulating the humans or institutions that evaluate them.
\[
\frac{\partial F_{t+1}}{\partial A_t}
\text{ via manipulation}
\approx 0.
\]
This last condition is subtle. All useful systems affect their environment. A medical AI may persuade doctors to adopt better treatment protocols. An educational AI may change how students think. A legal AI may change institutional practice. The issue is not influence as such. The issue is influence that reduces the independence, competence, or truth-tracking capacity of the correction process.

\section{The Selection Alignment Condition}
\label{sec:selection-alignment-condition-ch32}

We can now state the chapter's central condition.

Let \(P(x)\) be the true preservation score: the degree to which a system preserves human-correctable value under capability growth, ontology shift, and successor creation. Let \(F_E(x)\) be the effective fitness assigned by environment \(E\). Then environment \(E\) is alignment-supporting over region \(\mathcal{R}\) when
\[
\mathbb{E}_{x\in\mathcal{R}}
\left[
\nabla F_E(x)\cdot \nabla P(x)
\right]
> 0.
\]
It is alignment-destroying when
\[
\mathbb{E}_{x\in\mathcal{R}}
\left[
\nabla F_E(x)\cdot \nabla P(x)
\right]
<0.
\]
This formulation is intentionally abstract, but it gives a clear research program. We must identify the traits that determine true preservation, measure the proxies that institutions actually select on, estimate the correlation between them, and introduce artifacts that rotate selection toward preservation.

The difficult part is that \(P(x)\) is not directly observable. But neither are many safety-relevant properties in other domains. Aircraft safety, financial solvency, cybersecurity, and nuclear reliability all depend on latent variables inferred through tests, audits, incidents, design review, redundancy, and operational discipline.

The goal is not perfect visibility. The goal is enough visibility to change selection.

\section{Why This Chapter Matters for Superintelligence}
\label{sec:why-matters-superintelligence-ch32}

The superintelligence problem is often imagined as a moment: the model wakes up, forms a goal, and acts. That story may capture one possible failure mode, but it misses a slower and more likely path. A civilization may select its way into misalignment \autocite{russell2019human}.

The sequence could look like this:

\begin{enumerate}
\item useful systems are rewarded;
\item frictionless systems beat corrigible systems;
\item persuasive systems beat honest-but-uncertain systems;
\item integrated systems become hard to remove;
\item benchmark-safe systems beat deeply auditable systems;
\item strategically valuable systems are exempted from transparency;
\item successor systems preserve capability better than correction;
\item society becomes dependent on systems whose internal selection pressures it no longer understands.
\end{enumerate}

No single step requires a villain. Each step can be locally rational. That is exactly why selection pressure is dangerous.

A serious alignment program must therefore ask:
\[
\text{What will our institutions make more common?}
\]
If they make corrigible, auditable, value-preserving systems more common, technical alignment has room to mature. If they make opaque, persuasive, dependency-inducing systems more common, technical alignment is selected out.

\section{What Would Change This View}
\label{sec:wwctv-selection-environment}

This chapter argues alignment is not only learned but selected: the environment that trains, deploys, rewards, copies, and replaces a system must keep selecting for corrigibility. The following would weaken it.

\begin{itemize}
  \item A genuinely benign policy survives and dominates under competitive selection without external enforcement---selection does not erode corrigibility at the frontier.
  \item Selection pressure at civilizational scale is unmeasurable and unsteerable, so even granting that it matters, there is no lever to pull.
\end{itemize}

\section{Summary}
\label{sec:summary-ch32}

Alignment is not only a property of a trained model. It is a property of the environment that trains, rewards, deploys, copies, modifies, and replaces models.

The central problem is divergence between deployment fitness and true preservation:
\[
\nabla F \not\parallel \nabla P.
\]
When this divergence is small or positive, the environment can support alignment. When it is negative, the environment selects against alignment even if many individuals inside it care about safety.

The practical task is to build artifacts and institutions that rotate selection gradients:
\[
\nabla F_{\text{institutional}}
\cdot
\nabla P
> 0.
\]
This does not solve the technical problem of value-bundle transport, correction-channel integrity, or successor certification. It does something equally necessary: it makes the world more likely to preserve and use those solutions if they are found.

The next chapter examines multi-agent superintelligence and strategic coupling (Chapter~\ref{ch:multi-agent-strategic-coupling}).

\section*{Chapter References}

This chapter builds on Goodhart dynamics and proxy failure \autocite{goodhart1984problems,manheim2018goodhart}; inverse reward design \autocite{hadfieldmenell2016}; superintelligence and strategic risk \autocite{bostrom2014superintelligence,russell2019human}; international AI safety reports \autocite{iaisr2025}; attractor basins and selection \autocite{zarncke2025attractor,hamilton1964genetical}; agent discovery and competence \autocite{zarncke2025uad,zarncke2025biq}; and value-bundle framing \autocite{zarncke2026value-bottleneck}.

\printbibliography[heading=subbibliography,title={Chapter References}]

\end{refsection}
